\section{Introdução}
\label{sec:intro}

A última década consolidou o problema do \textit{alinhamento} de modelos de
linguagem como uma das frentes centrais da pesquisa em inteligência artificial
aplicada. Modelos pré-treinados em grandes \textit{corpora} de texto exibem
comportamentos emergentes desejáveis, mas também replicam padrões tóxicos,
estereótipos e violações de norma observados nos dados de
treinamento~\cite{bender2021,gehman2020}. A auditoria de tais modelos depende
criticamente da qualidade dos \textit{datasets} usados como instrumentos de
diagnóstico --- e, sobretudo, da \textit{representatividade} dos subconjuntos
sobre os quais inferências são realizadas.

A linhagem de \textit{datasets} rotulados para detecção de toxicidade tem
crescido ao longo dos últimos sete anos. \textit{Wikipedia Talk
Corpus}~\cite{wulczyn2017} e \textit{Jigsaw Toxic Comment
Classification}~\cite{borkan2019nuanced} estabeleceram a tarefa em larga
escala em inglês; \textit{HateXplain}~\cite{mathew2021} introduziu rótulos
explicativos; e, em português brasileiro, \textit{ToLD-Br}~\cite{leite2020} e
\textit{HateBR}~\cite{vargas2022hatebr} ampliaram a cobertura linguística e
cultural. Em comum, esses \textit{corpora} apresentam classes positivas com
desbalanceamento extremo, frequentemente abaixo de $1\%$ para categorias raras
como \textit{racism} ou \textit{threat}.

A teoria amostral clássica oferece os instrumentos necessários para esta
tarefa. Desenvolvida em ciências sociais e epidemiologia ao longo do
século~XX --- de Neyman~(1934) sobre alocação ótima a
Cochran~(1977)~\cite{cochran1977}, passando por
Bolfarine~\&~Bussab~(2005)~\cite{bolfarine2005} no contexto brasileiro, e
sintetizada em
Lohr~(1999)~\cite{lohr1999} e Thompson~(2002)~\cite{thompson2002} ---, esta
tradição formaliza derivação de tamanho amostral, alocação entre estratos e
estimadores que exploram variáveis auxiliares (razão e regressão) para
reduzir variância. Sua transposição para a auditoria de \textit{corpora} de
PLN é, contudo, recente e fragmentada.

Apesar do crescimento dos \textit{corpora}, a literatura empírica raramente
justifica estatisticamente o tamanho dos subconjuntos analisados.
\textit{Splits} pré-definidos ou amostras por conveniência (por exemplo, os
primeiros $N$ exemplos) são utilizados sem qualquer derivação formal de
margem de erro ou nível de confiança~\cite{leite2020,vargas2022hatebr}. Esta
prática, ainda que aceitável em treino preditivo, é insuficiente para o
propósito de \textit{auditoria}: estimar parâmetros populacionais (proporções,
médias, intensidades) sobre o \textit{corpus} requer teoria amostral
clássica~\cite{cochran1977,bolfarine2005,lohr1999,thompson2002}.

O presente trabalho propõe e implementa um \textit{pipeline} reproduzível de
auditoria estatística aplicado ao ToLD-Br, fundamentado em três decisões de
desenho:
\begin{enumerate}
  \item \textbf{Amostragem estratificada sem reposição}, com amostragem
        aleatória simples sem reposição (SRS-WOR) dentro de cada estrato.
  \item \textbf{Alocação ótima de Neyman}~\cite{cochran1977,bolfarine2005},
        que aloca esforço amostral em função do produto $N_h \sigma_h$ em cada
        estrato.
  \item \textbf{Estimadores tipo razão e tipo regressão}~\cite{cochran1977},
        em variantes \textit{separada} e \textit{combinada}, comparados ao
        \textit{baseline} da média estratificada simples.
\end{enumerate}

\subsection*{Questões orientadoras}

\begin{itemize}
  \item \textbf{Q1.} Qual o tamanho amostral mínimo --- pela fórmula de
        Cochran --- para estimar o parâmetro populacional médio com erro de
        amostragem de $5\%$ e confiança de $95\%$ sobre o ToLD-Br?
  \item \textbf{Q2.} Sob esse orçamento $n$, em que magnitude a alocação ótima
        de Neyman reduz o erro-padrão das estimativas frente à alocação
        proporcional clássica?
  \item \textbf{Q3.} Entre os estimadores tipo razão e tipo regressão (com a
        contagem de palavras como variável auxiliar), qual apresenta menor
        erro contra a verdade populacional, e o que isso revela sobre a
        relação linear entre toxicidade e indicadores léxicos no ToLD-Br?
\end{itemize}

\subsection*{Organização do texto}

A Seção~\ref{sec:theory} apresenta o referencial teórico de amostragem
estratificada, alocação de Neyman e estimadores tipo razão e regressão. A
Seção~\ref{sec:related} discute trabalhos relacionados em \textit{datasets} de
toxicidade e em aplicações de teoria amostral em PNL. A
Seção~\ref{sec:method} descreve a metodologia experimental: a população
ToLD-Br, a definição dos estratos, o cálculo de $n$ e a estratégia de
estimação. A Seção~\ref{sec:results} apresenta os resultados
empíricos. Por fim, a Seção~\ref{sec:conclusion} sintetiza as conclusões e
indica trabalhos futuros.
